% !TEX root = ../main.tex
\section{Metodología}

El sistema implementa el ciclo clásico de inferencia Mamdani, ilustrado en la figura~\ref{fig:flujo-mamdani}. Cada bloque corresponde a un módulo del paquete \texttt{fuzzy-core}, lo que garantiza una correspondencia uno a uno entre la teoría y la implementación.

\begin{figure}[H]
  \centering
  \begin{tikzpicture}[
      node distance=0.7cm and 0.5cm,
      every node/.style={font=\footnotesize,align=center},
      block/.style={rectangle,rounded corners=4pt,draw=difLine,line width=0.8pt,fill=difTeal!10,minimum width=2.4cm,minimum height=0.9cm,inner sep=4pt},
      result/.style={rectangle,rounded corners=4pt,draw=difLine,line width=0.8pt,fill=difCritical!12,minimum width=2.4cm,minimum height=0.9cm,inner sep=4pt},
      arrow/.style={-{Latex[length=2.5mm,width=1.8mm]},line width=0.8pt,difMuted},
    ]
    \node[block] (inputs) {Entradas crisp\\$(x_1,\ldots,x_5)$};
    \node[block,right=of inputs] (fuzz) {Fuzzificación\\$\mu_{A_i}(x_i)$};
    \node[block,right=of fuzz] (rules) {Reglas IF-THEN\\$\alpha_r=\min(\mu_{A_i})$};
    \node[block,below=of rules] (impl) {Implicación\\$\mu_{B'_r}=\min(\alpha_r,\mu_{B_r})$};
    \node[block,left=of impl] (agg) {Agregación\\$\mu_B=\max_r \mu_{B'_r}$};
    \node[block,left=of agg] (defuzz) {Defuzzificación\\$y^*$ por centroide};
    \node[result,below=of defuzz] (out) {Salida crisp\\y etiqueta};
    \draw[arrow] (inputs) -- (fuzz);
    \draw[arrow] (fuzz) -- (rules);
    \draw[arrow] (rules) -- (impl);
    \draw[arrow] (impl) -- (agg);
    \draw[arrow] (agg) -- (defuzz);
    \draw[arrow] (defuzz) -- (out);
  \end{tikzpicture}
  \caption{Flujo de inferencia Mamdani implementado. Cada bloque está respaldado por un módulo del motor \texttt{fuzzy-core} (\texttt{membership.ts}, \texttt{rules.ts}, \texttt{mamdani.ts}, \texttt{centroid.ts}).}
  \label{fig:flujo-mamdani}
\end{figure}

\subsection{Algoritmo formal}

\Needspace{16\baselineskip}
\begin{lstlisting}[caption={Pseudocódigo del motor difuso},label=lst:pseudocodigo]
entrada: promedio, asistencia, entregas, participacion, examenes
salida: riesgo academico crisp y etiqueta linguistica

para cada variable de entrada v:
  para cada conjunto difuso A en v.terminos:
    mu[v][A] = evaluar_pertenencia(A.forma, v.valor)

para cada regla r en base_de_reglas:
  alpha_r = min { mu[a.variable][a.termino] | a en r.antecedentes }
  para cada y en [0, 100] paso 1:
    salida_recortada_r(y) = min(alpha_r, mu_consecuente(y))

para cada y en [0, 100] paso 1:
  mu_agregada(y) = max { salida_recortada_r(y) | r en base }

numerador  = sum( y * mu_agregada(y) )
denominador = sum( mu_agregada(y) )
riesgo = numerador / denominador  si denominador > 0  sino 0
etiqueta = argmax_{B en salida.terminos} mu_B(riesgo)
\end{lstlisting}

\subsection{Discretización del universo de salida}

El universo de salida $[0,100]$ se discretiza con paso $\Delta y = 1$. Esta resolución es suficiente para estimar el centroide con un error relativo menor al $0.5\%$ comparado con la integral analítica, y mantiene la complejidad lineal $\mathcal{O}(101 \cdot |R|)$ donde $|R|$ es el número de reglas. Aumentar la resolución a $\Delta y = 0.1$ multiplica por diez el costo computacional sin cambios apreciables en la decisión final.

\subsection{Restricciones metodológicas}

El sistema es determinista y explicable. No utiliza entrenamiento, redes neuronales, regresión, estadística predictiva, APIs externas ni bases de datos. Toda decisión proviene de funciones de pertenencia y reglas visibles, lo que cumple los criterios de defensibilidad académica buscados en este proyecto. La trazabilidad completa se manifiesta en tres niveles:

\begin{itemize}
  \item \textbf{Estructural:} cada módulo del motor implementa una etapa única del proceso (figura~\ref{fig:flujo-mamdani}).
  \item \textbf{Numérica:} todas las figuras y tablas del reporte se generan a partir del motor mediante el script \texttt{emit-data}, sincronizado con cada compilación.
  \item \textbf{Operativa:} la interfaz web reproduce los valores, grados y reglas activadas de manera idéntica al pseudocódigo del listado~\ref{lst:pseudocodigo}.
\end{itemize}
